Modern Machine Learning (ML) often uses collaborative frameworks like Federated Learning (FL) to train models without collecting sensitive private data in a central server. Decentralized Federated Learning (DFL) takes this a step further by removing the central server entirely, relying instead on direct peer-to-peer communication. However, eliminating this central orchestrator creates a new problem: the network's physical layout becomes a huge part of the defense. This leaves the system highly vulnerable to targeted backdoor attacks, which are stealthy modifications where an attacker secretly teaches the model to misclassify specific trigger images without hurting its normal day-to-day performance. This thesis explores exactly how a network's shape helps or hinders the spread of these hidden attacks. 

We simulated a continuous attack across four different network structures: Fully Connected, Ring, Star, and a community-based Stochastic Block Model (SBM). To make the test fair and isolate the real effect of the topology, the attack had to survive the network's natural defense: the way honest nodes constantly average their models together to dilute anomalies. To achieve this, we introduced a specific formula that scales the attacker's malicious updates based on how many neighbors they connect to. This ensured the backdoor survived the averaging process without breaking the local model or immediately exposing the attacker. 

Testing this on both the MNIST and the industrial NEU-DET datasets proved that network geometry is a real bottleneck for an attacker. Dense, highly connected networks spread the infection rapidly, while sparse networks physically slow it down. Most importantly, realistic community-based networks naturally act as firewalls. They successfully trap malicious updates within local clusters, buying the wider network critical time before the infection can spread. Ultimately, this research shows that simply structuring a decentralized network correctly provides a powerful, built-in defense against AI poisoning attacks.
